\documentclass[11pt]{article}
\usepackage[T1]{fontenc}
\usepackage[utf8]{inputenc}
\usepackage{lmodern}
\usepackage[a4paper,margin=25mm]{geometry}
\usepackage{amsmath,amssymb,amsthm,mathtools}
\usepackage{microtype}
\usepackage{xcolor,hyperref}
\hypersetup{colorlinks=true,urlcolor=blue!55!black,linkcolor=blue!40!black,citecolor=blue!55!black}
\newtheorem{theorem}{Theorem}[section]
\newtheorem{proposition}[theorem]{Proposition}
\newtheorem{remark}[theorem]{Remark}
\newtheorem{definition}[theorem]{Definition}
\newcommand{\Q}{\mathbb{Q}}
\newcommand{\R}{\mathbb{R}}
\newcommand{\N}{\mathbb{N}}
\newcommand{\IT}{I_T}
\newcommand{\JT}{J_T}
\newcommand{\support}{\operatorname{supportT}}
\newcommand{\gram}{\operatorname{gramGen}}
\newcommand{\code}[1]{\texttt{#1}}
\title{Exact Rational Integration over Constrained Simplices\\[4pt]\large Symmetric Polynomial Bases, Marginal Geometry, and Lean 4 Verification}
\author{Reading Point \texorpdfstring{(\href{https://translatepi.ai/}{translatepi.ai})}{(translatepi.ai)}}
\date{Draft v0.6 --- 20 October 2026}
\begin{document}
\maketitle
\begin{abstract}
We present a formalized collection of integration identities for polynomial weights on
simplexes constrained by coordinate thresholds and cardinality-dependent caps.
The development includes exact simplex moments, a rational polynomial subspace of
 dimension 846, cap-parametric formulas for monomial integrals, rational Gram-entry
 representations, fibre decompositions, and two-aggregate Dirichlet integration.
  Exact polygon-cell formulas cover all marginal endpoint branches. Kernel-checked
  rational certificates establish the constant entry $J_T[1,1]$, the cross-entry
  $J_T[P_2,1]$, and the power-sum-square entry $J_T[P_2,P_2]$ under Table~3
  parameters. The V2h proof uses two distinguished-coordinate placements and
  a reusable rational moment table, completing a $2\times2$ marginal Gram block
  on $\operatorname{span}\{1,P_2\}$. The V2h Comparator's 27 statements were
  accepted by the Lean kernel on the VPS. The previously recorded V2g regression
  accepted 19 of 19 suites; no retained aggregate log certifies the full V2h
  regression count. These independent integration results contribute exact
  foundations for subsequent variational positivity work.
\end{abstract}

\section{Introduction and scope}
Exact integration on constrained regions is useful where a variational inequality must
ultimately be checked using rational arithmetic. A conventional simplex has a closed
monomial moment formula; the present support also depends on which coordinates cross a
threshold and on a cap for the sum of those coordinates. This introduces a finite
stratification, but preserves a route to exact, rational expressions.

We study the mathematical constructions in the public Lean repository:
\begin{center}
\href{https://github.com/thinkthoughts/lean-bgp212}{\nolinkurl{github.com/thinkthoughts/lean-bgp212}}
\end{center}
The motivating bounded-prime-gap argument originates with Axiom Math
\cite{axiom}; the statements below are an independent formalization of specified
components and do not amount to a complete formal proof of the unconditional
prime-gap bound. References to Lean declarations are included to delimit precisely
what has been checked. The short arguments given here are \emph{proof outlines}, not
substitutes for the complete Lean proof terms.

\section{Polynomial moments on a simplex}
\begin{definition}
For $m\in\N$ and $R\in\R$, define
\[
 \Delta_m(R)=\left\{x\in\R^m: x_i\ge 0\ \text{for all }i,
 \quad\sum_{i=1}^m x_i\le R\right\}.
\]
\end{definition}
\begin{theorem}[Simplex monomial moments]\label{thm:simplex}
For $m\in\N$, $R\ge0$, and $e\in\N^m$,
\[
 \int_{\Delta_m(R)}\prod_{i=1}^{m}x_i^{e_i}\,dx
 =\frac{R^{m+\sum_i e_i}\prod_i(e_i!)}{(m+\sum_i e_i)!}.
\]
Here the dimension-zero case uses the corresponding Lean integral convention.
\end{theorem}
\noindent\textbf{Lean:} \code{simplex\_integral}; definitions \code{simplex}, \code{momentR}.\\
\textit{Source:} \nolinkurl{Solution/SimplexMoment.lean:160}.

\paragraph{Proof outline.}
Induct on dimension and use Fubini to isolate a coordinate. The inner integral is a
Beta integral; the factorial expression then simplifies by induction. The formal
development additionally treats negative and zero radii, simplex volume, and
rational-polynomial integrals (\code{poly\_integral}).

\section{An indexed rational polynomial space}
Consider rational polynomials in 45 coordinates. Write $P_j$ for the $j$th power sum,
$P_{\lambda}$ for the product over parts of a multiset $\lambda$, and set
\[
 G_{a,\lambda}=(L-P_1)^aP_{\lambda},\qquad
 \lambda\text{ has positive even parts},\quad a+|\lambda|\le21.
\]
The index type $B_{21}\mathrm{Idx}$ enumerates the admissible pairs.
\begin{theorem}[Independence and rank]\label{thm:basis}
The indexed family $G_{a,\lambda}$ is linearly independent over $\Q$. Its rational span
$B_{21}$ has
\[
 \dim_{\Q}B_{21}=846.
\]
\end{theorem}
\noindent\textbf{Lean:} \code{genIdx\_linearIndependent},
\code{finrank\_B21} (\nolinkurl{Solution/Variational.lean:398,410}).

\paragraph{Proof outline.}
The indexing is injective; distinct indexed monomials are linearly independent.
Injective polynomial transformations preserve this independence. A separate count
establishes $|B_{21}\mathrm{Idx}|=846$, and the span-rank identity gives the result.
The theorem concerns this designated subspace, not the full polynomial ring.

\section{Cap-constrained supports and general integrals}
Let $N\in\N$, $g\ge0$, $L\in\R$, and $C:\N\to\R$. For $x\in\R^N$ write
\[
 R_g(x)=\{i:x_i>g\},\qquad s_g(x)=|R_g(x)|,\qquad
 \rho_g(x)=\sum_{i\in R_g(x)}x_i.
\]
Define the capped simplex support
\begin{equation}\label{eq:support}
 T(N,g,L,C)=\left\{x\in\R^N:x_i\ge0,\quad
 \sum_i x_i\le L,\quad \rho_g(x)\le C(s_g(x))\right\}.
\end{equation}
\begin{theorem}[General monomial integral]\label{thm:general}
For $N\in\N$, $g\ge0$, $L\in\R$, $C:\N\to\R$, $e\in\N^N$, and $p\in\N$,
\begin{equation}\label{eq:gram}
 \int_{T(N,g,L,C)}\left(L-\sum_i x_i\right)^p
 \prod_i x_i^{e_i}\,dx
 =\gram(N,g,L,C,e,p).
\end{equation}
The right-hand side is the explicit, finite expression defined in Lean.
\end{theorem}
\noindent\textbf{Lean:} \code{integral\_supportT\_gen}; \code{gramGen}.\\
\textit{Sources:} \nolinkurl{Solution/GramGeneral.lean:344};\\
\nolinkurl{BGP212/GramGeneral.lean}.

\paragraph{Proof outline.}
Stratify by the rough-coordinate set, apply inclusion--exclusion to threshold
conditions, and translate the resulting pieces to weighted simplexes. Use
Theorem~\ref{thm:simplex} and Fubini to evaluate the polynomial weights.
Placement-counting identities \code{sig\_insert}, \code{sig\_union}, and
\code{sig\_card} reduce the sums to the finite closed form. The reduction
\code{gramGen\_extend} transports calculation to the coordinates having nonzero
exponents. This is an algebraic efficiency result, not a claim that all entries are
computationally cheap to certify.

\section{Rational Gram representations}
Specialize to the fixed rational cap values and support parameters of Table 3 in
\cite{axiom}, with $N=45$. For indices $i,j\in B_{21}\mathrm{Idx}$ define
\[
 \IT(i,j)=\int_T G_i(x)G_j(x)\,dx.
\]
\begin{theorem}[Explicit rational representation]\label{thm:rational}
For every $i,j\in B_{21}\mathrm{Idx}$, $\IT(i,j)$ equals the real embedding of a
finite rational sum indexed by placements $\varphi$ of the combined power-sum parts:
\[
 \IT(i,j)=\left(\sum_{\varphi}
 \gram_{\Q}\bigl(45,g,L,C,e_{\varphi},a_i+a_j\bigr)\right)_{\R}.
\]
Here $e_{\varphi}$ is the exponent vector obtained from the placement, exactly as
specified by \code{expo} in the Lean source.
\end{theorem}
\noindent\textbf{Lean:} \code{gram\_entry\_B21\_rat}
(\nolinkurl{Solution/GramGeneral.lean:522});
\code{gram\_entry\_gen\_list} and \code{gram\_entry\_example}.

The theorem gives each entry an explicit rational representation; it does not
supply an evaluated and certified $846\times846$ matrix. One worked entry,
corresponding to $(L-P_1)P_2^2$, has its value certified using a collection of
small exact rational lemmas. This separation between an algebraic identity and
its numerical certificate is important to reproducibility.

\section{Marginal-fibre geometry}
Let $\widehat x\in\R^n$ and append $u\in\R$ as the last coordinate. Define
$S=\sum_i\widehat x_i$, $s=s_g(\widehat x)$,
$\rho_1=\rho_g(\widehat x)$, and
\[
 \widehat T=\{\widehat x\ge0:S\le\tau,\ \rho_1\le C(s)\},
 \quad U(\widehat x)=\min\!\left\{L-S,\max\bigl(g,C(s+1)-\rho_1\bigr)\right\}.
\]
\begin{theorem}[Fibre equivalence]\label{thm:fibre}
Suppose $g\ge0$, $\tau\le L$, and $\widehat x\in\widehat T$. Then for every
$u\in\R$,
\[
 (\widehat x,u)\in T(n+1,g,L,C)
 \quad\Longleftrightarrow\quad 0\le u\le U(\widehat x).
\]
\end{theorem}
\noindent\textbf{Lean:} \code{fibre\_equiv}
(\nolinkurl{Solution/Fibre.lean:115}).

\paragraph{Proof outline.}
Appending $u$ leaves the rough set unchanged for $u\le g$, and increments both
rough count and rough mass for $u>g$. Combining the resulting cap inequality with
$S+u\le L$ identifies the minimum--maximum endpoint exactly.
The empty-fibre theorem for cap-violating outer coordinates additionally needs
$C(m+1)\le C(m)+g$ for $m\ge1$, and $C(0)\ge0$. No constraint on the step
$m=0$ is assumed; the relevant Table 3 facts are checked in Lean.

\begin{theorem}[Marginal identity]\label{thm:marginal}
Assume $g\ge0$, $\tau\le L$, $C(0)\ge0$, and
$C(m+1)\le C(m)+g$ for every $m\ge1$. For continuous $P$,
\[
 \operatorname{marginalJ}(n,g,L,\tau,C,P)
 =\JT(n,g,L,\tau,C,P).
\]
The right-hand functional integrates the squared fibre integral over $\widehat T$.
\end{theorem}
\noindent\textbf{Lean:} \code{marginal\_agrees}
(\nolinkurl{Solution/Fibre.lean:217}).

\begin{proposition}[Polynomial inner integral]\label{prop:inner}
For arbitrary $U,M\in\R$ and $p,t\in\N$,
\[
 \int_0^U(M-u)^pu^t\,du
 =\sum_{j=0}^p\binom pj M^{p-j}(-1)^j
 \frac{U^{t+j+1}}{t+j+1}.
\]
\end{proposition}
\noindent\textbf{Lean:} \code{inner\_integral}
(\nolinkurl{Solution/Fibre.lean:240}). The equality follows by expanding the
binomial and integrating each power. It holds for arbitrary real $U$ under the
oriented interval-integral convention.

\section{Two-aggregate Dirichlet integration}\label{sec:aggregate}
Let $I$ and $J$ be disjoint, nonempty finite coordinate groups, with exponent
vectors $a\in\N^I$ and $b\in\N^J$. For a nonempty finite group $G$ and an
exponent vector $e$, define
\[
 \alpha_G(e)=|G|+\sum_{i\in G}e_i-1,\qquad
 A_G(e)=\frac{\prod_{i\in G}(e_i!)}{\alpha_G(e)!}.
\]
Write $\Delta_{I\sqcup J}(M)$ for the nonnegative simplex with total coordinate
sum at most $M$, and
\[
 \Delta_2(M)=\{(u_1,u_2)\in\R_{\ge0}^2:u_1+u_2\le M\}.
\]
\begin{theorem}[Two-aggregate Dirichlet integration]\label{thm:aggregate}
Suppose $I$ and $J$ are nonempty, and $h:\R^2\to\R$ is measurable and
bounded on $\Delta_2(M)$. With $u_1=\sum_{i\in I}y_i$ and
$u_2=\sum_{j\in J}y_j$, the weighted simplex integral satisfies
\begin{align}\label{eq:twoagg}
 &\int_{\Delta_{I\sqcup J}(M)}
   \left(\prod_{i\in I}y_i^{a_i}\right)
   \left(\prod_{j\in J}y_j^{b_j}\right)
   h(u_1,u_2)\,dy \notag\\
 &\qquad=A_I(a)A_J(b)
   \int_{\Delta_2(M)} u_1^{\alpha_I(a)}u_2^{\alpha_J(b)}
   h(u_1,u_2)\,du_1\,du_2.
\end{align}
The formalization also gives the version on $\mathrm{Fin}\,N$ for a rough
coordinate set and its complement.
\end{theorem}
\noindent\textbf{Lean:} \code{two\_group}, \code{two\_group\_fin},
\code{one\_group}; \nolinkurl{Solution/Aggregate.lean}. The theorem is proved
for the corresponding measurable, bounded integrands; separate theorems handle
empty groups and other boundary cases.

\paragraph{Proof outline.}
Use Fubini and a sum-coordinate shear to isolate the total in one group;
apply the weighted simplex moment calculation to the corresponding slices.
Repeat for the second group. The factors $A_I,A_J$ record the Dirichlet
normalizations; in particular the $-1$ in each aggregate exponent is essential.
The polynomial cases \code{tri\_poly} and \code{two\_group\_poly} yield
explicit rational expressions, and \code{two\_group\_const} recovers the
simplex moment formula when $h=1$.

\begin{proposition}[A cut-triangle certificate]\label{prop:aggregate-example}
With $y=(y_0,y_1,y_2)$, the exact value is
\[
 \int_{\substack{y_i\ge0,\;y_0+y_1+y_2\le1\\y_0\le 1/2}}
 y_0y_2^2\,dy=\frac{19}{7680}.
\]
\end{proposition}
\noindent\textbf{Lean:} \code{tri\_cut}, \code{example\_value}
(\nolinkurl{Solution/Aggregate.lean}). The certificate reduces this integral
to $\tfrac13\int_{\Delta_2(1),\,u_1\le1/2}u_1u_2^3\,du_1du_2$;
the value has also been cross-checked by direct rational integration.

\paragraph{Connection to marginal geometry.}
On the translated stratum $x=y+g\mathbf 1_{R\cup K}$, define
$u_1=\sum_{i\in R}y_i$ and $u_2=\sum_{i\notin R}y_i$.
The proved correspondence \code{stratumEndpoint\_shift} expresses the
fibre endpoint as
\[
 U=\min\!\left(L-(r+k)g-u_1-u_2,\;
 \max\bigl(g,C(r+1)-rg-u_1\bigr)\right),
 \qquad r=|R|,\;k=|K|.
\]
The cap index belongs to the stratum $R$. Substituting the translated point's
rough set $R\cup K$ instead changes the endpoint; the Table~3 control records
the distinct values $587/1250$ and $461/1250$.
The cap-branch polygon-cell integration is established in Section~\ref{sec:capcell}; all endpoint branches and the constant marginal Gram entry are treated in Section~\ref{sec:v2f}; the power-sum entry follows in Section~\ref{sec:v2g}, and the power-sum-square entry in Section~\ref{sec:v2h}.


\section{Polygon-cell integration and the cap branch}\label{sec:capcell}
The two-aggregate reduction produces planar regions bounded by coordinate
thresholds and a total-sum constraint. For real $A,B,D$ define the clipped cell
\[
 \mathcal B(A,B,D)=\{(u_1,u_2)\in\R^2:0\le u_1\le A,\;
 0\le u_2\le B,\;u_1+u_2\le D\}.
\]
\begin{theorem}[Exact polygon-cell moments]\label{thm:box}
For $\alpha,\beta\in\N$, the Lebesgue integral
\[
 \int_{\mathcal B(A,B,D)}u_1^\alpha u_2^\beta\,du_1\,du_2
 =\operatorname{boxVal}(A,B,D,\alpha,\beta)
\]
where $\operatorname{boxVal}$ is a finite closed form using polynomial
antiderivatives. For rational $A,B,D$ its value is rational. The same
construction gives rational integrals for arbitrary rational polynomials in
two variables. Degenerate and empty cells are treated separately.
\end{theorem}
\noindent\textbf{Lean:} \code{box\_monomial}, \code{boxVal\_cast},
\code{box\_poly}, \code{boxVal\_empty}, \code{boxVal\_nonpos},
\code{boxVal\_rect}; \nolinkurl{Solution/CapCell.lean}.

\paragraph{Proof outline.}
Integrate first in $u_2$ with upper endpoint
$\min(B,D-u_1)$ and split the $u_1$ interval at $D-B$. The resulting
one-variable powers have exact polynomial antiderivatives. In particular,
when $A+B\le D$ the diagonal constraint is redundant and the cell is a
rectangle; the formula also covers diagonally truncated cells.

\begin{theorem}[Cap-branch endpoint and piece integral]\label{thm:capbranch}
On a translated stratum $x=y+g\mathbf1_{R\cup K}$, with
$r=|R|$, $k=|K|$ and aggregates $u_1,u_2$ as in
Section~\ref{sec:aggregate}, the cap branch of the fibre endpoint is
\[
 U=C(r+1)-rg-u_1,
\]
subject to its vertical and horizontal branch constraints. Combined with
the piece constraints, the aggregate domain is a clipped cell
$\mathcal B(A,B,D)$ with parameters determined by $g,L,\tau,C,r,k$.
The weighted piece integrals
\[
 \int_{\mathrm{piece}} y^t\,\mathbf1_{\mathrm{cap}}
 U^m\,dy
\]
therefore equal an explicit expression obtained from the two-aggregate
Dirichlet constants and rational cell integral formulas.
\end{theorem}
\noindent\textbf{Lean:} \code{endpointAgg\_on\_capCell},
\code{cap\_branch\_piece}, \code{cap\_branch\_r0},
\code{cap\_inner}; \nolinkurl{Solution/CapCell.lean}.
The formalization proves the branch classification and corresponding
endpoint equalities. For Table 3 the $\tau$ constraint is redundant on
the cap branch: $\tau-(L-g)=1/625>0$, and the cap cells become rectangles
(\code{tau\_ge\_table3}, \code{capCell\_rect\_table3}).

\begin{proposition}[A Table 3 cap-piece certificate]\label{prop:capcert}
For the specified 44-coordinate piece $(R,K)=(\{0\},\varnothing)$,
the cap-branch contribution to the integral defining $J_T(1)$ is
\[
 \frac{(L-C(2))^{43}\big((C(2)-g)^3-g^3\big)}{3\cdot43!}
 \approx 1.0464\times 10^{-70}.
\]
\end{proposition}
\noindent\textbf{Lean:} \code{capPieceT3\_value}
(\nolinkurl{Solution/CapCell.lean}). This is \emph{one piece} of the
inclusion--exclusion decomposition, not the complete value of $J_T(1)$.
The exact rational certificate is kernel-checked. Polygon clipping and
lower-dimensional integration provide independent numerical controls;
the latter is an approximation, not part of the kernel proof.
The remaining branches and their assembly are treated in Section~\ref{sec:v2f}.

\section{All endpoint branches and the first marginal Gram entry}\label{sec:v2f}
The translated marginal endpoint is the minimum of the remaining total mass
and the maximum of a cap expression and $g$:
\[
 U=\min\bigl(L-(r+k)g-u_1-u_2,\;
       \max(g,C(r+1)-rg-u_1)\bigr).
\]
Its switching lines partition each two-aggregate piece into four cells:
the cap branch, the remaining-mass branch on the cap side, the constant
$g$ branch, and the remaining-mass branch on the $g$ side. Each cell
is represented by signed differences of the clipped polygon cells
$\mathcal B(A,B,D)$ in Section~\ref{sec:capcell}.

\begin{theorem}[Exact branch-cell and piece integrals]\label{thm:allbranches}
For nonnegative integer exponents, the weighted integral on each branch cell
of a polynomial block
\[
 u_1^\alpha u_2^\beta (p-u_1)^{m_1}
 (L'-u_1-u_2)^{m_2}
\]
reduces to a finite combination of exact clipped-cell moments. The
corresponding integrals on translated rough-set strata admit closed forms,
including the stratum with no rough coordinates. For integrands depending
on the stratum, inclusion--exclusion and cardinality counting assemble
these pieces into the marginal functional.
\end{theorem}
\noindent\textbf{Lean:} \code{box\_twoPow}, \code{cells\_integral},
\code{piece\_block}, \code{piece\_r0},
\code{indicator\_supportT\_strat}, \code{count\_by\_card};
\nolinkurl{Solution/MarginalEntry.lean}.

\begin{theorem}[Rational reduction of the constant-generator marginal]\label{thm:jtone}
For $N\ge1$ and parameters satisfying $0\le g\le C(1)$ and
$C(N)<Ng$, the squared-fibre marginal for the constant polynomial is
expressed by the explicit rational-piece expression $\operatorname{jtSum}$:
\[
 J_T(1)=(N+1)\int_{\widehat T}U(\widehat x)^2\,d\widehat x
 =\operatorname{jtSum}(N,g,L,\tau,C,2,0).
\]
\end{theorem}
\noindent\textbf{Lean:} \code{jT\_one\_eq},
\code{jT\_one\_pieces}, \code{piece\_eq\_pieceJ};
\nolinkurl{Solution/MarginalEntry.lean}. The reduction uses binomial
counting for this symmetric integrand; entries with power-sum factors
require additional expansion and placement accounting.

\begin{proposition}[Table 3 marginal Gram certificate]\label{prop:jtcert}
Under the Table 3 parameters, for the generator $G_{0,\varnothing}=1$,
\[
 J_T[G_{0,\varnothing},G_{0,\varnothing}]
 =45\int_{\widehat T}U(\widehat x)^2\,d\widehat x
 \approx 1.1193487646397893\times10^{-55}.
\]
The exact value is the ratio of these integers, obtained by concatenating each
integer's digit lines (no separators):
\begin{center}
\small
\begin{tabular}{@{}l@{}}
\textbf{Numerator}\\
\ttfamily
196405301342598493830210603911544843582810204\\
587055080769623409436245894353192939488221872\\
912132381844932553462651053065161718578698716\\
871038141\\[4pt]
\textbf{Denominator}\\
\ttfamily
175463901463993193504011955433551437477604239\\
482178113383469748717206470461915657210738572\\
757691144943237304687500000000000000000000000\\
000000000000000000000000000000000000000000000\\
0000000000000
\end{tabular}
\end{center}
\end{proposition}
\noindent\textbf{Lean:} \code{jtSum\_table3},
\code{jT\_one\_table3}, \code{eval\_gen\_zero},
\code{jT\_gram\_gen00\_table3};
\nolinkurl{Solution/MarginalEntryCert.lean},
\nolinkurl{Solution/MarginalEntry.lean}.
The proof evaluates 133 nonempty $(r,k)$ pieces by kernel checking and
proves the remaining pieces vanish. The exact certificate is regenerated
byte-for-byte by \code{marginal\_entry\_v2f\_emit.py}.
Independent polygon clipping, low-dimensional integration, and Monte Carlo
provide additional checks, but are not substitutes for Lean verification.
This establishes the constant-generator entry. The nonconstant power-sum
entry is developed in Section~\ref{sec:v2g}.


\section{Bilinear marginal form and a power-sum Gram entry (V2g)}\label{sec:v2g}
Write $P_2(x)=\sum_i x_i^2$, and let $\widehat P_2$ denote the same
power sum on the first $44$ coordinates. The bilinear marginal functional
$J_T[P,Q]$ polarizes the squared-fibre integral. Its symmetric bilinearity
is formalized by \code{jTbil\_self} and \code{jTbil\_comm}.

\begin{proposition}[Fibre polynomial for $P_2$]\label{prop:p2inner}
For a fibre of length $U$ above $\widehat x$, the constant and quadratic
polynomial fibre integrals satisfy
\[
 \int_0^U 1\,du=U,\qquad
 \int_0^U P_2(\widehat x,u)\,du
 =U\widehat P_2(\widehat x)+\frac{U^3}{3}.
\]
Consequently the cross-entry is
\begin{equation}\label{eq:p2cross}
 J_T[P_2,1]
 =45\int_{\widehat T}\left(U^2\widehat P_2+\frac{U^4}{3}\right)d\widehat x.
\end{equation}
\end{proposition}
\noindent\textbf{Lean:} \code{inner\_P2}, \code{inner\_one},
\code{p2Entry\_eq}; \nolinkurl{Solution/PowerSumEntry.lean}.

\paragraph{Coordinate placements.}
On a translated stratum $x=y+g\mathbf1_{R\cup K}$, the power sum becomes
\[
 \widehat P_2(x)=\sum_i y_i^2
   +2g\sum_{i\in R\cup K}y_i+g^2|R\cup K|.
\]
One distinguished coordinate lies in $R$, in $K$, or in neither;
these cases carry multiplicities $r$, $k$, and $N-r-k$.
The resulting exponent shifts feed the exact two-aggregate Dirichlet
weights and polygon-cell formulas established earlier.
Theorems \code{piece\_single}, \code{piece\_single\_eq},
\code{sum\_place}, and \code{piece\_P2} establish the piecewise and
cardinality-weighted assembly, including the necessary boundary cases.

\begin{proposition}[Table 3 nonconstant marginal certificate]\label{prop:p2cert}
Under the Table 3 parameters, the marginal Gram cross-entry has an exact
rational value $q_{P_2,1}\in\Q$, certified in Lean, with
\[
 J_T[P_2,1]=q_{P_2,1}
 \approx 4.999439687639965\times10^{-57}.
\]
The complete numerator and denominator are preserved by the Lean
certificate and the reproducible certificate generator.
\end{proposition}
\noindent\textbf{Lean:} \code{p2EntryVal\_table3},
\code{jT\_P2\_one\_table3}, \code{jT\_one\_P2\_table3}.\par
\noindent\textit{Sources:} \nolinkurl{Solution/PowerSumEntry.lean},
\nolinkurl{Solution/PowerSumEntryCert.lean}.

The exact Table 3 evaluation uses 133 cardinality-class pieces with two
certificate lemmas per piece (266 lemmas total); the remaining classes
are eliminated by \code{cond\_tail} and \code{cond\_high}.
The generator \code{specs/checks/power\_sum\_entry\_v2g\_emit.py}
reproduces the Lean certificate byte-for-byte on the VPS.
Six independent Lean-file controls have their expected outcomes: two
accepted positive controls and four rejected negative controls.
Numerical cross-checks include a low-dimensional grid and Monte Carlo;
these support the arithmetic independently of the exact kernel proof.
The power-sum-square diagonal entry $J_T[P_2,P_2]$, including two
distinguished-coordinate placements and their diagonal case, is proved in Section~\ref{sec:v2h}.

\section{A two-coordinate moment table and the power-sum-square entry (V2h)}\label{sec:v2h}
The marginal square of the nonconstant generator is
\begin{align}\label{eq:p2square}
 J_T[P_2,P_2]
 &=45\int_{\widehat T}\left(U\widehat P_2+\frac{U^3}{3}\right)^2d\widehat x\\
 &=45\int_{\widehat T}\left(U^2\widehat P_2^2
   +\frac23 U^4\widehat P_2+\frac19 U^6\right)d\widehat x.
\end{align}
The new term contains two distinguished coordinate indices:
\[
 \widehat P_2^2=\sum_{i,j}x_i^2x_j^2.
\]
For $i=j$ a translated summand becomes $(y_i+s_i)^4$; for $i\ne j$
we obtain $(y_i+s_i)^2(y_j+s_j)^2$, where $s_i$ is the stratum
translation ($g$ for coordinates in $R\cup K$, zero otherwise).
The three placement classes $R$, $K$, and their complement have sizes
$n_c$; the number of ordered distinct-coordinate pairs in classes
$c,d$ is $n_cn_d-\mathbf1_{c=d}n_c$. The diagonal contributes $n_c$
placements per class.

\begin{theorem}[Moment-table reduction for the marginal square]\label{thm:v2hmoment}
For each translated stratum and each exponent vector $t$, the
weighted piece integral of $y^t U^m$ reduces to a factorial factor
$\prod_i t_i!$ multiplied by an exact moment determined by the total
monomial degrees in the rough and smooth coordinate groups. Together
with the diagonal and distinct-coordinate placement formulas, this
gives an exact finite rational representation of~\eqref{eq:p2square}.
\end{theorem}
\noindent\textbf{Lean development:}
\nolinkurl{BGP212/PowerSumPlace.lean},\newline\nolinkurl{BGP212/PowerSumSquare.lean},
\nolinkurl{Solution/PowerSumSquare.lean};\newline
\nolinkurl{Challenge/PowerSumSquare.lean}.
The proof reuses the two-aggregate and polygon-cell machinery, with
moments shared among coordinate placements rather than recomputed for
each summand.

\begin{proposition}[Table 3 power-sum-square certificate]\label{prop:p2squarecert}
At the specified Table~3 parameters, the entry $J_T[P_2,P_2]$ is an
exact rational number represented by the V2h Lean certificate, with
\[
 J_T[P_2,P_2]\approx 2.27107\times10^{-58}.
\]
The certificate is organized into 14 resumable row files and a final
assembly file, each checked as part of the V2h Comparator proof.
\end{proposition}
\noindent\textbf{Certificate sources:}
\nolinkurl{Solution/PowerSumSquareCert0.lean} through\newline
\nolinkurl{Solution/PowerSumSquareCert13.lean},\newline
\nolinkurl{Solution/PowerSumSquareCertSum.lean};\newline
\nolinkurl{specs/checks/power_sq_entry_v2h_emit.py}.
The 27-statement Comparator configuration is
\nolinkurl{Comparator/power_sum_square_v2h.json}.
Its final VPS output states that the Lean default kernel accepts the
solution and reports ``Your solution is okay!''.

\paragraph{The first $2\times2$ marginal Gram block.}
Combining V2f, V2g, and V2h yields the exact rational entries of
\[
 \mathcal G_{\{1,P_2\}}=
 \begin{pmatrix}
  J_T[1,1] & J_T[1,P_2]\\
  J_T[P_2,1] & J_T[P_2,P_2]
 \end{pmatrix}
 \approx
 \begin{pmatrix}
  1.11935\times10^{-55} & 4.99944\times10^{-57}\\
  4.99944\times10^{-57} & 2.27107\times10^{-58}
 \end{pmatrix}.
\]
Using the displayed approximations, the Cauchy--Schwarz ratio
\[
 \frac{J_T[P_2,1]^2}{J_T[1,1]\,J_T[P_2,P_2]}\approx0.984
\]
is close to one. This numerical observation motivates exact assessment
of determinants and variational certificates; an independent Lean
certificate for a strictly positive $2\times2$ determinant is not
asserted here. Independent Monte Carlo and low-dimensional grid
comparisons test the computed value but are distinct from the Lean proof.

\section{Verification, limitations, and further mathematics}
The project uses Lean \code{v4.34.0-rc2} and a pinned Mathlib revision
\code{85e3a25}. At the V2g baseline, all 19 Comparator suites were independently
accepted on the VPS. Four positive V2d controls were accepted; four intentionally
incorrect controls were rejected with unsolved goals. The V2d proof commit is
\code{52fd36a}; its published verification record is commit
\href{https://github.com/thinkthoughts/lean-bgp212/commit/90055db}{\code{90055db}}.
For V2e-A, five positive controls were accepted and four negative controls
were rejected. The V2e-A proof commit is \code{01644e6}, with VPS
verification/provenance commit
\href{https://github.com/thinkthoughts/lean-bgp212/commit/38606c0}{\code{38606c0}}.
For V2e-B, five positive controls were accepted and five deliberately
incorrect controls rejected; the V2e-B proof commit is
\href{https://github.com/thinkthoughts/lean-bgp212/commit/6a21780}{\code{6a21780}},
with verification/provenance commit
\href{https://github.com/thinkthoughts/lean-bgp212/commit/25a7376}{\code{25a7376}}.
For V2f, four positive Lean-file controls were accepted and five negative
controls rejected. Its proof commit is
\href{https://github.com/thinkthoughts/lean-bgp212/commit/8c50fb6}{\code{8c50fb6}},
with VPS provenance commit
\href{https://github.com/thinkthoughts/lean-bgp212/commit/da0c3fd}{\code{da0c3fd}}.
For V2g, two positive controls were accepted and four negative controls
rejected; its 266 certificate lemmas regenerated byte-identically on the
VPS. The V2g proof commit is
\href{https://github.com/thinkthoughts/lean-bgp212/commit/ad70abb}{\code{ad70abb}},
with verification/provenance commit
\href{https://github.com/thinkthoughts/lean-bgp212/commit/12f8b16}{\code{12f8b16}}.
For V2h, the 27-statement Comparator accepted the power-sum-square
certificate on the VPS. Its proof commit is
\href{https://github.com/thinkthoughts/lean-bgp212/commit/6cef1f4}{\code{6cef1f4}},
with VPS provenance commit
\href{https://github.com/thinkthoughts/lean-bgp212/commit/a9b0e0c}{\code{a9b0e0c}}.
The recorded V2h proof includes 14 rational moment-certificate rows and
an assembly theorem. The final Comparator acceptance is confirmed; a
complete 20/20 regression count has not been established from a saved log.
The reported kernel axioms for these developments are
\code{propext}, \code{Quot.sound}, and \code{Classical.choice}.
Results and verification scripts are available in
\href{https://github.com/thinkthoughts/lean-bgp212/tree/main/specs/results}{\code{specs/results}}.

Further work extends the three certified entries to general
power-sum marginal Gram entries and their exact rational evaluation,
optimized polynomial coefficients, exact variational positivity certificates, and
remaining analytic hypotheses. The existing Lean development formalizes the
\emph{conditional} implication $\mathrm{DHL}[45,2]\Rightarrow H_1\le212$;
this manuscript does not claim an unconditional formal proof of $H_1\le212$.

\section*{Acknowledgment and provenance}
The bounded-prime-gap application is motivated by Axiom Math's mathematical argument
\cite{axiom}. This manuscript records an independent Lean formalization of the
specified results, with theorem names and commit identifiers to support reproducible
review. Substantive exposition and attribution will be expanded in a subsequent draft.

\begin{thebibliography}{9}
\bibitem{axiom}
Axiom Math, \emph{A New Bound for Small Gaps Between Primes}, preliminary draft,
3 September 2026. Project page:
\url{https://primegaps.axiommath.ai/}.
\bibitem{leanrepo}
Reading Point, \emph{lean-bgp212}: Lean source, verification records, and certificates.
\url{https://github.com/thinkthoughts/lean-bgp212}.
\end{thebibliography}
\end{document}
